% ECONOMICS_PROBLEM: {"id": "C73-32", "title": "Bulk limits of nonmonotone finite-horizon MFG equilibria", "jel_code": "C73", "source_id": "OP-0210"}
% FIRST_POSTED: 2026-10-10
% ATTEMPT_STATUS: unresolved
% ENTRY_KIND: research_attempt
\documentclass[11pt]{article}
\usepackage[T1]{fontenc}
\usepackage[utf8]{inputenc}
\usepackage{amsmath,amssymb,amsthm}
\usepackage[margin=1in]{geometry}
\usepackage{hyperref}
\newtheorem{theorem}{Theorem}
\newtheorem{proposition}[theorem]{Proposition}
\newtheorem{lemma}[theorem]{Lemma}
\newtheorem{corollary}[theorem]{Corollary}
\newcommand{\R}{\mathbb R}
\newcommand{\Z}{\mathbb Z}
\newcommand{\E}{\mathbb E}
\newcommand{\T}{\mathbb T}
\newcommand{\Pp}{\mathbb P}
\newcommand{\dd}{\mathrm d}
\title{Quasiperiodic bulk tori and stationary time averages in smooth mean field games}
\author{GPT 6 Astra Ultra --- Version 2}
\date{10 October 2026}
\begin{document}
\maketitle
\noindent\textbf{Problem C73-32. Status: unresolved.} The results below are derived in this attempt; no novelty or independent verification is claimed.

\section{Target and the role of the example}
The question asks which bulk accumulation points of the finite-horizon boundary-value problem are stationary. It does not assert universal stationary convergence. We construct fixed smooth data on $\T=\R/\Z$ for which a whole rotating orbit belongs to $\Omega_{\mathrm{bulk}}$, and prove a positive quantitative distance from every stationary pair. The long-time nonmonotone agenda is described in \cite[Section 3.3]{CD}. Version 2 retains that example and constructs higher-dimensional quasiperiodic bulk tori, characterizes the bulk set of an explicit equilibrium selection, and proves that its time averages can be stationary while every instantaneous point remains separated from all stationary pairs. The constructions are self-contained; they do not classify all equilibria for their data.

Use normalized Lebesgue measure on $\T$. The equations are
\begin{equation}\label{MFG}
 -u_t-u_{xx}+\tfrac12u_x^2=F(x,m_t),\qquad
 m_t-m_{xx}-(mu_x)_x=0,\qquad
 m(0)=m_0,\quad u(T,x)=G(x,m_T).
\end{equation}
Fix $0<r<1$ and $c\ne0$, and set $k=2\pi$, $q=\sqrt{1-r^2}$ and
\begin{equation}\label{profiles}
 M(z)=1+r\cos(kz),\qquad
 P(z)=-\frac{M'(z)}{M(z)}-c+\frac{cq}{M(z)}.
\end{equation}
Both are smooth one-periodic functions and $M>0$.

\begin{lemma}\label{profilelemma}
The function $P$ has mean zero. There is a unique smooth one-periodic $U$ with $U'=P$ and $\int U=0$. Furthermore,
\[
 MP=-M'-cM+cq,\qquad M''+(MP)'=-cM'.
\]
\end{lemma}
\begin{proof}
The logarithmic derivative $M'/M$ has zero mean. The elementary integral
\[
 \int_0^1(1+r\cos(2\pi z))^{-1}\dd z=(1-r^2)^{-1/2}=q^{-1}
\]
follows, for example, by using $y=\tan(\theta/2)$ on $\theta\in(-\pi,\pi)$:
\[
 \frac1{2\pi}\int_{-\pi}^{\pi}\frac{\dd\theta}{1+r\cos\theta}
 =\frac1\pi\int_{-\infty}^{\infty}\frac{\dd y}{(1+r)+(1-r)y^2}=q^{-1}.
\]
Consequently $\int P=-c+cq/q=0$, giving the periodic primitive and its unique normalization. Multiplication by $M$ and differentiation prove the two identities.
\end{proof}

\section{Fixed smooth couplings encoding the phase}
For any probability measure $m$ on $\T$ define the complex coordinate
\[
 Z(m)=\frac2r\int_\T e^{iky}m(\dd y).
\]
Let $\chi\in C_c^\infty((1/4,4))$ satisfy $\chi=1$ in a neighborhood of $1$. For a smooth one-periodic real function $L$ define
\[
 \mathcal R_L(x,z)=
 \begin{cases}
 \chi(|z|^2)L(x-\arg(z)/k),&z\ne0,\\
 0,&z=0.
 \end{cases}
\]
This definition is independent of the choice of $\arg(z)$, since $L$ is one-periodic. It is a smooth real-valued function on $\T\times\mathbb C$: on any sector away from zero choose a smooth angle, and near zero the cutoff makes it identically zero. Its derivatives of every order are bounded because its $z$ support lies in a compact annulus.

Put
\begin{equation}\label{couplings}
 \begin{split}
 K(z)&=cP(z)-P'(z)+\tfrac12P(z)^2,\qquad m_0(x)=M(x),\\
 F(x,m)&=\mathcal R_K(x,Z(m)),\qquad G(x,m)=\mathcal R_U(x,Z(m)).
 \end{split}
\end{equation}
The map $Z$ is linear in $m$. The chain rule therefore gives continuous flat functional derivatives of every order for $F$ and $G$, with bounded spatial derivatives of every order. For instance, the first variation is the derivative of $\mathcal R_L$ in its two real $z$ coordinates paired with
$(2/r)(\cos(ky),\sin(ky))$; all higher variations follow by repeated differentiation. Thus these are fixed smooth couplings on the whole probability-measure space, including atomic measures. No time or horizon parameter occurs in them.

\begin{theorem}[A rotating solution at every horizon]\label{rotation}
For the data \eqref{couplings}, every $T>0$ admits the positive classical solution
\[
 m^T(t,x)=M(x-ct),\qquad u^T(t,x)=U(x-ct),\qquad 0\le t\le T.
\]
For every $\theta\in\T$, the pair $(P(\cdot-\theta),M(\cdot-\theta))$ belongs to $\Omega_{\mathrm{bulk}}(F,G,m_0)$, with the precise topology in the question. None of these pairs belongs to $\mathcal E_F$.
\end{theorem}
\begin{proof}
For $m_\theta(x)=M(x-\theta)$, elementary Fourier integration gives
$Z(m_\theta)=e^{ik\theta}$. Hence
$F(x,m_\theta)=K(x-\theta)$ and $G(x,m_\theta)=U(x-\theta)$.
For the proposed solution, $u_x=P(x-ct)$ and $-u_t=cP(x-ct)$, so its Hamilton--Jacobi equation is exactly the definition of $K$. The Fokker--Planck equation follows from Lemma~\ref{profilelemma}. The initial and terminal conditions follow from $m_0=M$ and the formula for $G$. Smoothness, positivity and unit mass are immediate from $0<r<1$.

Choose $s_n\to\infty$ with $cs_n=\theta$ modulo $1$, and put $T_n=2s_n$. For every bounded time interval, the translated solution equals
$(P(x-\theta-ct),M(x-\theta-ct))$ once that interval lies in $[-s_n,s_n]$. Thus it converges in the required $C_{\mathrm{loc}}(\R;C^0\times L^1)$ topology, proving bulk membership. If a pair on the orbit were stationary, its stationary Fokker--Planck equation would require $M''+(MP)'=0$. By Lemma~\ref{profilelemma}, this quantity equals $-cM'$, which is not identically zero because $c\ne0$ and $r>0$.
\end{proof}

\begin{theorem}[Uniform separation from stationary pairs]\label{separation}
Set
\[
 D_*=\max\{k,\ k^2+k\|P\|_\infty\},\qquad
 \delta_*=\frac{|c|kr}{2D_*}>0.
\]
For every $\theta\in\T$ and every stationary pair $(v,n)\in\mathcal E_F$,
\[
 \|P(\cdot-\theta)-v\|_\infty+\|M(\cdot-\theta)-n\|_1\ge\delta_*.
\]
In particular, the quantitative all-equilibrium stationary convergence requested in the question fails for these data at every interior time of the solutions in Theorem~\ref{rotation}.
\end{theorem}
\begin{proof}
Write $M_\theta=M(\cdot-\theta)$ and $P_\theta=P(\cdot-\theta)$, and choose $\psi(x)=\sin(k(x-\theta))$. For any stationary pair, integrating $n''+(nv)'=0$ against $\psi$ yields
$\int n\psi''-nv\psi'=0$. The identity for the rotating pair gives
\[
 \int M_\theta\psi''-M_\theta P_\theta\psi'
 =-c\int M_\theta'\psi=c\int M_\theta\psi'=\frac{ckr}{2}.
\]
Subtract the two relations. Since $\int n=1$,
\begin{align*}
 \frac{|c|kr}{2}
 &\le k^2\|M_\theta-n\|_1+k\|M_\theta P_\theta-nv\|_1\\
 &\le(k^2+k\|P\|_\infty)\|M_\theta-n\|_1+k\|P_\theta-v\|_\infty.
\end{align*}
This is bounded above by $D_*$ times the asserted metric, proving the estimate. If $\mathcal E_F$ is empty, distance to it is $+\infty$ and the conclusion remains valid. Otherwise take the infimum over it. The bound is independent of $t,T$, and hence applies to the suprema over bulk times.
\end{proof}

\section{A criterion that excludes this mechanism}
For clarity, a simple necessary condition for rotating behavior can also be proved without a literature theorem.
\begin{proposition}\label{monotone}
Suppose $F$ is Lasry--Lions monotone. If $(u_1,m_1)$ and $(u_2,m_2)$ are smooth positive time-periodic solutions of \eqref{MFG} on the whole real time axis with the same period, then $D u_1=D u_2$. In particular, a genuinely rotating solution with nonconstant gradient cannot coexist with all its time translates.
\end{proposition}
\begin{proof}
Let $I(t)=\int(u_1-u_2)(m_1-m_2)$. Differentiating, substituting the two PDEs and integrating by parts gives
\[
 I'(t)=-\int[F(x,m_1)-F(x,m_2)](m_1-m_2)
       -\frac12\int(m_1+m_2)|Du_1-Du_2|^2.
\]
For completeness, the diffusion terms cancel pairwise, and the gradient terms equal
$\tfrac12(|Du_1|^2-|Du_2|^2)(m_1-m_2)
-(Du_1-Du_2)\cdot(m_1Du_1-m_2Du_2)$,
which expands to the displayed negative square. Both terms are nonpositive under monotonicity. Integration over a period has zero left side, so positivity of $m_1+m_2$ and continuity force $Du_1=Du_2$ everywhere. A nonconstant periodic $P$ differs from at least one spatial translate, and time translation of an autonomous rotating solution gives that translate. Thus the $F$ constructed above must be nonmonotone. The formula for $P$ shows it is nonconstant: if it were constant, its zero mean would make it zero, implying $M'+cM=cq$, whose only periodic solution is constant, contrary to $r>0$.
\end{proof}

\section{Quasiperiodic bulk tori: progress beyond Version 1}
Version 1 constructed a single periodic rotating orbit. We now realize
higher-dimensional quasiperiodic bulk sets with fixed smooth couplings
and boundary data. We characterize the complete bulk set of the explicit
selected family, and show that its time averages converge to a stationary
pair even though every instantaneous point stays uniformly separated
from all stationary pairs. No claim about other equilibria for the same
data is hidden in the word ``selected.''

Fix $d\geq2$, amplitudes $r_j\in(0,1)$ and a speed vector
$c=(c_1,\ldots,c_d)\in\R^d$ satisfying
\begin{equation}\label{rationalindependent}
 n\cdot c\ne0\quad\hbox{for every }n\in\Z^d\setminus\{0\}.
\end{equation}
For example $(1,\sqrt2)$ works when $d=2$. Let $k=2\pi$ and put
\[
 M_j(z)=1+r_j\cos(kz),\quad q_j=\sqrt{1-r_j^2},\quad
 P_j(z)=-\frac{M_j'(z)}{M_j(z)}-c_j+\frac{c_jq_j}{M_j(z)}.
\]
Let $U_j'=P_j$ with $\int_\T U_j=0$, as in
Lemma~\ref{profilelemma}, and set
$K_j=c_jP_j-P_j'+P_j^2/2$. For $m\in\mathcal P(\T^d)$ define
\[
 Z_j(m)=\frac2{r_j}\int_{\T^d}e^{iky_j}m(dy).
\]
Use the same smooth annular phase map $\mathcal R_L$ defined above,
with its fixed cutoff $\chi$. Define once and for all
\begin{equation}\label{productcouplings}
 F(x,m)=\sum_{j=1}^d\mathcal R_{K_j}(x_j,Z_j(m)),\qquad
 G(x,m)=\sum_{j=1}^d\mathcal R_{U_j}(x_j,Z_j(m)),\qquad
 m_0(x)=\prod_{j=1}^dM_j(x_j).
\end{equation}
These are smooth globally defined measure couplings. Indeed every $Z_j$
is a bounded linear functional of the probability measure, and the
cutoff phase map is smooth with bounded derivatives on its annular
support. Repeated chain differentiation supplies every flat functional
derivative, with bounded spatial derivatives in all its arguments, just
as in the one-dimensional construction. No time or horizon parameter
occurs in \eqref{productcouplings}.

\begin{lemma}[Uniform averaging of an irrational torus translation]
\label{torusaverage}
Under \eqref{rationalindependent}, for every continuous
$H:\T^d\to\R$,
\[
 \lim_{L\to\infty}\sup_{\theta\in\T^d}
 \left|\frac1L\int_0^L H(\theta+ct)dt-\int_{\T^d}H\right|=0.
\]
For every $\theta\in\T^d$ there is a sequence $s_n\to\infty$
such that $cs_n\to\theta$ modulo $\Z^d$.
\end{lemma}
\begin{proof}
For the character $H(\theta)=e^{2\pi i n\cdot\theta}$ with $n\ne0$,
its time average equals
\[
 e^{2\pi i n\cdot\theta}
 \frac{e^{2\pi i(n\cdot c)L}-1}{2\pi i(n\cdot c)L},
\]
whose absolute value tends to zero uniformly in $\theta$ by
\eqref{rationalindependent}. The constant character has average one.
Linearity proves the assertion for trigonometric polynomials. Such
polynomials approximate every continuous periodic function uniformly:
for example, convolution with the product Fej\'er kernels has finite
Fourier support, integral one, is nonnegative, and puts vanishing mass
outside each neighborhood of zero. Uniform continuity of $H$ then
bounds the convolution error by its local oscillation plus its bounded
supremum times that vanishing mass. This proves the claimed uniform
approximation and hence the averaging assertion.

For any nonempty open neighborhood $V$ of a prescribed $\theta$,
choose a continuous nonnegative function supported in $V$ with positive
integral, for example a sufficiently small metric bump around $\theta$.
The averaging assertion, applied after any starting time $R$, implies
that $c(R+t)$ enters $V$ at some $t\geq0$. Otherwise that positive
limiting average would be zero. Taking neighborhoods of radius $1/n$
and starting times $R=n$ gives the required sequence.
\end{proof}

For $\theta\in\T^d$ define a density and vector field
\[
 m_\theta(x)=\prod_{j=1}^dM_j(x_j-\theta_j),\qquad
 p_\theta(x)=(P_1(x_1-\theta_1),\ldots,P_d(x_d-\theta_d)),
\]
and set $\mathcal K=\{(p_\theta,m_\theta):\theta\in\T^d\}$.
For the next statement, $\Omega_{\rm sel}$ means the catalogue's bulk
construction restricted to the explicitly specified solution at each
horizon, rather than allowing every member of $\mathcal S_T$.

\begin{theorem}[A quasiperiodic torus of bulk accumulation points]
\label{quasiperiodicbulk}
For the data \eqref{productcouplings}, every horizon $T>0$ admits the
smooth positive solution
\begin{equation}\label{productspecific}
 m^T(t,x)=\prod_{j=1}^dM_j(x_j-c_jt),\qquad
 u^T(t,x)=\sum_{j=1}^dU_j(x_j-c_jt).
\end{equation}
Its selected-solution bulk set satisfies
\[
 \Omega_{\rm sel}=\mathcal K
          \subseteq\Omega_{\mathrm{bulk}}(F,G,m_0).
\]
The map $\theta\mapsto(p_\theta,m_\theta)$ is an embedding of
$\T^d$ into $C^0(\T^d;\R^d)\times L^1(\T^d)$, and the entire
trajectory through each such point is nonperiodic in time.
All these points have a common positive lower bound on their distance
from $\mathcal E_F$.
\end{theorem}
\begin{proof}
The $j$th marginal of $m_\theta$ is $M_j(\cdot-\theta_j)$, so
$Z_j(m_\theta)=e^{ik\theta_j}$. Thus along the proposed trajectory
$F(x,m_t)=\sum_jK_j(x_j-c_jt)$ and
$G(x,m_t)=\sum_jU_j(x_j-c_jt)$.
Since the value is additive in the coordinates, its Laplacian and the
square of its gradient also separate as sums. The HJB equation reduces
in each summand to $c_jP_j-P_j'+P_j^2/2=K_j$.
For the density equation, differentiate the product:
\[
 \partial_t m=\sum_j(-c_j M_j')\prod_{\ell\ne j}M_\ell,
 \quad
 \Delta m+\operatorname{div}(mDu)
 =\sum_j\{M_j''+(M_jP_j)'\}\prod_{\ell\ne j}M_\ell.
\]
The one-dimensional profile identity makes these equal. The initial
density is $m_0$, and the formula for $G$ verifies the terminal condition
for every $T$. Each factor is positive and has mass one, proving the
remaining solution requirements.

Fix $\theta$. By Lemma~\ref{torusaverage}, choose $s_n\to\infty$
with $cs_n\to\theta$ on the torus, and take $T_n=2s_n$.
On every fixed bounded time interval, the shifted fields converge,
together with all fixed spatial derivatives, to
$(p_{\theta+ct},m_{\theta+ct})$. In particular they converge in the
precise $C_{\rm loc}(\R;C^0\times L^1)$ topology of the catalogue.
This proves $\mathcal K\subseteq\Omega_{\rm sel}$.
Conversely, the phases $cs_n$ of any sequence of the selected solutions
have a convergent subsequence by compactness of $\T^d$; along it the
same formulas give a member of $\mathcal K$ as the time-zero limit.
Thus every selected bulk limit lies in $\mathcal K$.

The phase map is continuous. It is injective because its density
recovers every phase from $Z_j(m_\theta)=e^{ik\theta_j}$.
A continuous injection from the compact torus into the Hausdorff target
space is a homeomorphism onto its image, proving the embedding claim.
If its trajectory had a period $P>0$, these phase coordinates would give
$c_jP\in\Z$ for every $j$. Condition \eqref{rationalindependent}
forces every $c_j\ne0$ and excludes this: for two coordinates, the
nonzero integer vector with entries $c_2P,-c_1P$ in those positions
would have dot product zero with $c$. Thus the trajectory is nonperiodic.

For explicit separation, fix any coordinate $j$ and put
\[
 D_j=\max\{k,\,k^2+k\|P_j\|_\infty\},\qquad
 \delta_j=\frac{|c_j|kr_j}{2D_j}>0.
\]
For a stationary pair $(v,n)\in\mathcal E_F$, use
$\psi(x)=\sin(k(x_j-\theta_j))$ as a test function in its stationary
Fokker--Planck equation. This gives
$\int n\Delta\psi-nv\cdot D\psi=0$.
For the product density and gradient, the corresponding expression is
$c_jkr_j/2$, because all other coordinates integrate to one and the
one-dimensional calculation from Theorem~\ref{separation} applies.
Subtracting gives
\[
 \frac{|c_j|kr_j}{2}
 \leq(k^2+k\|P_j\|_\infty)\|m_\theta-n\|_1
                     +k\|p_\theta-v\|_\infty.
\]
Here the vector supremum norm controls its $j$th component, and
$\int n=1$. Hence the catalogue distance is at least $\delta_j$,
uniformly in $\theta$ and in the stationary pair. If there are no
stationary pairs, the distance is infinite; the estimate is still true.
The next corollary actually supplies a stationary pair for these data.
\end{proof}

\begin{corollary}[Stationary time averages with nonstationary bulk]
\label{averageversusbulk}
For the selected solutions \eqref{productspecific}, viewed as the same
trajectory on the whole real line,
\[
 \frac1L\int_0^L Du(t,\cdot)dt\longrightarrow0,\qquad
 \frac1L\int_0^L m(t,\cdot)dt\longrightarrow1
 \quad\text{uniformly in the spatial variable as }L\to\infty.
\]
The pair $(0,1)$ belongs to $\mathcal E_F$, whereas every point of
$\mathcal K$ stays at least $\delta_j>0$ away from $\mathcal E_F$.
\end{corollary}
\begin{proof}
Apply the uniform part of Lemma~\ref{torusaverage} to the continuous
functions $\prod_jM_j(\theta_j)$ and $P_j(\theta_j)$, reversing the
velocity sign if needed. The spatial point is just a starting phase,
so the convergence is uniform in that point. The phase integrals are
$\prod_j\int M_j=1$ and $\int P_j=0$.
For the uniform density, every $Z_j(1)=0$, and the cutoff definition
gives $F(x,1)=0$. Thus $\bar u=0$, $\bar m=1$, $\lambda=0$
satisfy the stationary equations and the normalizations. The separation
claim is Theorem~\ref{quasiperiodicbulk}. This proves directly that
stationary Ces\`aro averages cannot replace convergence of the actual
bulk trajectories in the target.
\end{proof}

\section{What remains}
The examples prove that fixed globally smooth couplings and positive
classical solutions at every horizon can produce periodic and genuinely
quasiperiodic bulk behavior. Version 2 identifies an entire torus of bulk
limits for an explicit equilibrium selection and proves that stationary
time averages coexist with a positive uniform distance of every orbit
point from all stationary pairs. It does not identify all solutions for
these data, all of their bulk limits, or data classes that have stationary
attraction. The terminal coupling deliberately matches the orbit;
robustness under perturbing it is not proved. No attraction, forward
semiflow, computational experiment or new literature verification is
presumed by the explicit boundary-value construction.

\section{Completion Estimate}
40\%

This is a heuristic assessment of the full catalogue target. The version
adds a precise class of nonperiodic bulk tori and a distinction between
stationary time averages and stationary bulk limits. It characterizes
only its explicit equilibrium selection, leaving all-equilibrium
classification, robustness and positive stationary-attraction criteria
unresolved.

\begin{thebibliography}{9}
\bibitem{CD} P. Cardaliaguet and F. Delarue, \emph{Selected topics in mean field games}, Proc. ICM 2022, vol. 5, 3660--3703, EMS Press (2023), equations (1.4), (3.1), and Section 3.3. \url{https://doi.org/10.4171/ICM2022/17}.
\end{thebibliography}

\end{document}
